\documentclass[10pt,a4paper]{amsart}
\usepackage{amsmath,amssymb,amsthm,graphicx}
% Retained minimal preamble for the Stacks Project chapter more-morphisms.tex (More on Morphisms).
% Verbatim excerpt of the chapter preamble: Xy-pic setup (native preamble lines 14-18) and the native
% theorem declarations with the mathematical macros (native preamble lines 154-217). Copied mechanically;
% nothing added, removed or redefined. The wrapper class and the amsmath/amssymb/amsthm/graphicx base
% packages are supplied by the assembler, so the class/externaldocument part of the native preamble is omitted.
% For commutative diagrams we use Xy-pic
\usepackage[all]{xy}

% We use 2cell for 2-commutative diagrams.
\xyoption{2cell}
% Theorem environments.
%
\theoremstyle{plain}
\newtheorem{theorem}[subsection]{Theorem}
\newtheorem{proposition}[subsection]{Proposition}
\newtheorem{lemma}[subsection]{Lemma}

\theoremstyle{definition}
\newtheorem{definition}[subsection]{Definition}
\newtheorem{example}[subsection]{Example}
\newtheorem{exercise}[subsection]{Exercise}
\newtheorem{situation}[subsection]{Situation}

\theoremstyle{remark}
\newtheorem{remark}[subsection]{Remark}
\newtheorem{remarks}[subsection]{Remarks}

\numberwithin{equation}{subsection}

% Macros
%
\def\lim{\mathop{\mathrm{lim}}\nolimits}
\def\colim{\mathop{\mathrm{colim}}\nolimits}
\def\Spec{\mathop{\mathrm{Spec}}}
\def\Hom{\mathop{\mathrm{Hom}}\nolimits}
\def\Ext{\mathop{\mathrm{Ext}}\nolimits}
\def\SheafHom{\mathop{\mathcal{H}\!\mathit{om}}\nolimits}
\def\SheafExt{\mathop{\mathcal{E}\!\mathit{xt}}\nolimits}
\def\Sch{\mathit{Sch}}
\def\Mor{\mathop{\mathrm{Mor}}\nolimits}
\def\Ob{\mathop{\mathrm{Ob}}\nolimits}
\def\Sh{\mathop{\mathit{Sh}}\nolimits}
\def\NL{\mathop{N\!L}\nolimits}
\def\CH{\mathop{\mathrm{CH}}\nolimits}
\def\proetale{{pro\text{-}\acute{e}tale}}
\def\etale{{\acute{e}tale}}
\def\QCoh{\mathit{QCoh}}
\def\Ker{\mathop{\mathrm{Ker}}}
\def\Im{\mathop{\mathrm{Im}}}
\def\Coker{\mathop{\mathrm{Coker}}}
\def\Coim{\mathop{\mathrm{Coim}}}

% Boxtimes
%
\DeclareMathSymbol{\boxtimes}{\mathbin}{AMSa}{"02}

%
% Macros for moduli stacks/spaces
%
\def\QCohstack{\mathcal{QC}\!\mathit{oh}}
\def\Cohstack{\mathcal{C}\!\mathit{oh}}
\def\Spacesstack{\mathcal{S}\!\mathit{paces}}
\def\Quotfunctor{\mathrm{Quot}}
\def\Hilbfunctor{\mathrm{Hilb}}
\def\Curvesstack{\mathcal{C}\!\mathit{urves}}
\def\Polarizedstack{\mathcal{P}\!\mathit{olarized}}
\def\Complexesstack{\mathcal{C}\!\mathit{omplexes}}
% \Pic is the operator that assigns to X its picard group, usage \Pic(X)
% \Picardstack_{X/B} denotes the Picard stack of X over B
% \Picardfunctor_{X/B} denotes the Picard functor of X over B
\def\Pic{\mathop{\mathrm{Pic}}\nolimits}
\def\Picardstack{\mathcal{P}\!\mathit{ic}}
\def\Picardfunctor{\mathrm{Pic}}
\def\Deformationcategory{\mathcal{D}\!\mathit{ef}}
\usepackage{cleveref}
\crefname{originallemma}{original lemma}{original lemmas}
\Crefname{originallemma}{Original lemma}{Original lemmas}
\makeatletter
\newlabel{lemma-check-perfect-stalks}{{069C}{source}{}{sourcepointer-1}{}}
\newlabel{lemma-check-perfect-stalks@cref}{{[originallemma][0][]069C}{[1][0][]source}{}{}{}}
\makeatother
\crefname{originallemma}{original lemma}{original lemmas}
\Crefname{originallemma}{Original lemma}{Original lemmas}
\makeatletter
\newlabel{lemma-pseudo-coherent-fppf-local-source}{{0699}{source}{}{sourcepointer-2}{}}
\newlabel{lemma-pseudo-coherent-fppf-local-source@cref}{{[originallemma][0][]0699}{[1][0][]source}{}{}{}}
\makeatother
\crefname{originalreference}{original reference}{original references}
\Crefname{originalreference}{Original reference}{Original references}
\makeatletter
\newlabel{more-algebra-lemma-flat-descent-tor-amplitude}{{068S}{source}{}{sourcepointer-3}{}}
\newlabel{more-algebra-lemma-flat-descent-tor-amplitude@cref}{{[originalreference][0][]068S}{[1][0][]source}{}{}{}}
\makeatother
\title{Perfectness of a morphism is fppf local on the source}
\author{The Stacks Project Authors}
\date{stacks-project:more-morphisms.tex:local-snapshot-20261011; GNU FDL 1.2 or later}
\begin{document}
\maketitle
\noindent\textit{Editorial note.} Counted span is the author's own proof environment, source lines 17997-18013, immediately following the statement at 17991-17995.\par\medskip
\noindent\textit{Editorial note.} Three cross-chapter pointers (Descent, More on Algebra) plus one same-chapter lemma are recorded in assembly.external\_references.\par\medskip
\noindent\textit{Editorial note.} The statement is a single-property locality claim; the two halves of perfectness are handled separately inside the span.\par\medskip
\noindent\textit{Allowed standard prerequisites (editorial).}\begin{itemize}
\item Pseudo-coherence is fppf local on the source. chapter lemma lemma-pseudo-coherent-fppf-local-source, cited in the proof. Supplies the pseudo-coherence half of the fppf locality argument.
\item Flatness of a flat locally finitely presented base change. Descent, lemma descent-lemma-flat-finitely-presented-permanence, cited in the proof. Keeps flatness along the fppf coverings that appear in the reduction.
\item Tor dimension modulo a nilpotent and its base change behaviour. More on Algebra lemmas on tor amplitude cited in the proof. Converts the stalkwise tor amplitude condition under the covering into the condition on X.
\end{itemize}\medskip
\ifdefined\hypertarget\hypertarget{sourcepointer-1}{}\fi
\noindent\textit{Source reference.} Original lemma 069C: sources/stacks-more-morphisms/pinned-original-source.tex, source lines 17901--17901. Navigation pointer only: the counted proof cites this statement by \textbackslash{}ref as a standing input; the deduction the counted proof has to make is carried out inside the counted span, and no part of that deduction is delegated to this pointer.\par\medskip
\ifdefined\hypertarget\hypertarget{sourcepointer-2}{}\fi
\noindent\textit{Source reference.} Original lemma 0699: sources/stacks-more-morphisms/pinned-original-source.tex, source lines 17621--17621. Navigation pointer only: the counted proof cites this statement by \textbackslash{}ref as a standing input; the deduction the counted proof has to make is carried out inside the counted span, and no part of that deduction is delegated to this pointer.\par\medskip
\ifdefined\hypertarget\hypertarget{sourcepointer-3}{}\fi
\noindent\textit{Source reference.} Original reference 068S: focused-seeds/geometry/S9-morph-perfect-fppf-local-source/ref-tags.txt, source lines 4--4. Navigation pointer only: the counted proof cites this statement by \textbackslash{}ref as a standing input; the deduction the counted proof has to make is carried out inside the counted span, and no part of that deduction is delegated to this pointer.\par\medskip
\begin{lemma}
\label{lemma-perfect-fppf-local-source}
The property $\mathcal{P}(f) =$``$f$ is perfect''
is fppf local on the source.
\end{lemma}

\begin{proof}
Let $\{g_i : X_i \to X\}_{i \in I}$ be an fppf covering of schemes and let
$f : X \to S$ be a morphism such that each $f \circ g_i$ is
perfect. By
Lemma \ref{lemma-pseudo-coherent-fppf-local-source}
we conclude that $f$ is pseudo-coherent.
Hence by
Lemma \ref{lemma-check-perfect-stalks}
it suffices to check that $\mathcal{O}_{X, x}$ is an
$\mathcal{O}_{S, f(x)}$-module of finite tor dimension for all $x \in X$.
Pick $i \in I$ and $x_i \in X_i$ mapping to $x$. Then we see that
$\mathcal{O}_{X_i, x_i}$ has finite tor dimension over
$\mathcal{O}_{S, f(x)}$ and that
$\mathcal{O}_{X, x} \to \mathcal{O}_{X_i, x_i}$ is faithfully flat.
The desired conclusion follows from
More on Algebra, Lemma \ref{more-algebra-lemma-flat-descent-tor-amplitude}.
\end{proof}

\end{document}
